% GENERATED FILE. Edit canonical lessons, then run npm run generate:books.
% canonical-source-hash: fnv1a64:16ea7e2d7d59d770
% canonical-lessons: GU-W04-lla, GU-W04-tha, GU-W04-independent-a, GU-W04-i-matra, GU-R04-first-four-r1

\chapter{Four Forms Before Your Name}
\label{ch:gu-first-use-forms-a}
\hlchaptermodality{4}

\begin{chapteropening}
\textbf{By the end of this chapter:} \emph{I can read and write four Gujarati forms before they appear in the first name exchange: \gu{ળ}, \gu{થ}, \gu{અ}, and \gu{િ}.}

With every model covered, read \gu{ળ}, \gu{થ}, \gu{અ}, and \gu{િ} and write each one from a spoken sound or role cue, keeping reading and writing scores separate.
\end{chapteropening}

\section[ḷa]{\gu{ળ} — one new shape, retroflex lla}

\label{lesson:GU-W04-lla}



\begin{quote}
Write retroflex \emph{nna} once from memory. Keep the tongue-curl idea, but put the old shape away: today it receives a different sound and form.
\end{quote}

\subsection*{Script — meet one form}
Keep \textbf{\gu{ળ}} visible. It is retroflex \emph{lla}: the tongue curls back for the \emph{l} sound. Begin at the upper left and circle the broad left bowl. Continue upward through the narrow middle turn, cross the high right arch, and descend the tall spine into its lower foot without lifting.

\subsection*{Writing — notice and trace}
Finger-trace \textbf{\gu{ળ}} once, then air-trace it once. Notice broad bowl, narrow turn, high arch, tall spine; no pen lift.

\subsection*{Writing — copy with the model}
Copy visible \textbf{\gu{ળ}} once. Compare only the broad bowl and tall spine; repair one landmark if needed.

\subsection*{Writing — recall after hiding}
Hide the model and write retroflex \emph{lla} once. Uncover \textbf{\gu{ળ}}, compare, and stop after one repair.

\begin{tcolorbox}[breakable,colback=teal!4,colframe=teal!35!black,title={Before you move on}]
Name the sound, the two main landmarks, and the pen-lift count.

Sources: \href{https://www.t30apps.com/gujarati-alphabet-writing/}{Gujarati Alphabet Writing Practice} and the \href{https://www.unicode.org/charts/PDF/U0A80.pdf}{Unicode Gujarati chart}.
\end{tcolorbox}
\section[tha]{\gu{થ} — the consonant tha}

\label{lesson:GU-W04-tha}



\begin{quote}
Write retroflex \emph{lla} once from memory. Check only its broad bowl and tall spine.
\end{quote}

\subsection*{Script — meet one form}
Keep \textbf{\gu{થ}} visible. It is \emph{tha}: say \emph{ta} with a clear puff of air. Begin at the right of the small upper loop, circle it, descend through the middle, and sweep around the broad lower body into the right shoulder. Lift once; descend the tall right spine and turn through its lower foot.

\subsection*{Writing — notice and trace}
Finger-trace \textbf{\gu{થ}} once, then air-trace it once. Notice small upper loop, broad lower body, separate right spine; one pen lift.

\subsection*{Writing — copy with the model}
Copy visible \textbf{\gu{થ}} once. Compare only the upper loop and separate spine; repair one landmark if needed.

\subsection*{Writing — recall after hiding}
Hide the model and write \emph{tha} once. Uncover \textbf{\gu{થ}}, compare, and stop after one repair.

\begin{tcolorbox}[breakable,colback=teal!4,colframe=teal!35!black,title={Before you move on}]
Say \emph{tha} with its puff, then name the pen-lift count.

Sources: \href{https://www.t30apps.com/gujarati-alphabet-writing/}{Gujarati Alphabet Writing Practice} and the \href{https://www.unicode.org/charts/PDF/U0A80.pdf}{Unicode Gujarati chart}.
\end{tcolorbox}
\section[a]{\gu{અ} — short a standing alone}

\label{lesson:GU-W04-independent-a}



\begin{quote}
Write \emph{tha} once from memory. Then say the vowel already carried by a bare Gujarati consonant: short \emph{a}.
\end{quote}

\subsection*{Script — meet one form}
Keep \textbf{\gu{અ}} visible. It is independent short \emph{a}, used when the vowel begins a syllable without a consonant. Sweep clockwise around the open left curve, continue through the lower body, rise into the middle shoulder, and finish the small right arch. Lift once; descend the separate right stem into its foot.

\subsection*{Writing — notice and trace}
Finger-trace \textbf{\gu{અ}} once, then air-trace it once. Notice open left curve, middle shoulder, small arch, separate right stem; one pen lift.

\subsection*{Writing — copy with the model}
Copy visible \textbf{\gu{અ}} once. Compare only the open curve and separate stem; repair one landmark if needed.

\subsection*{Writing — recall after hiding}
Hide the model and write independent short \emph{a} once. Uncover \textbf{\gu{અ}}, compare, and stop after one repair.

\begin{tcolorbox}[breakable,colback=teal!4,colframe=teal!35!black,title={Before you move on}]
When does short \emph{a} need \textbf{\gu{અ}} instead of arriving free on a consonant? (\textbf{When it stands without a consonant.})

Sources: \href{https://www.t30apps.com/gujarati-alphabet-writing/}{Gujarati Alphabet Writing Practice} and the \href{https://www.unicode.org/charts/PDF/U0A80.pdf}{Unicode Gujarati chart}.
\end{tcolorbox}
\section[-i]{\gu{િ} — short i sits before, sounds after}

\label{lesson:GU-W04-i-matra}



\begin{quote}
Point to \textbf{\gu{અ}} and \textbf{\gu{શ}}. Say independent short \emph{a}, then \emph{sha}. Do not combine them; today one small sign changes only the vowel on \textbf{\gu{શ}}.
\end{quote}

\subsection*{Script — meet one form}
Keep \textbf{\gu{િ}} visible. It is the attached short-\emph{i} sign. It is written to the left of its consonant even though \emph{i} is pronounced after that consonant. With \textbf{\gu{શ}}, it makes \textbf{\gu{શિ}}, \emph{shi}. Start with the upright stem and finish through the curl that reaches toward the consonant.

\subsection*{Writing — notice and trace}
Finger-trace \textbf{\gu{િ}} once, then air-trace it once. Notice the sign's left-side place in \textbf{\gu{શિ}}. Say \emph{shi} only after the whole form is visible.

\subsection*{Writing — copy with the model}
Copy visible \textbf{\gu{શિ}} once. Circle only the short-\emph{i} sign and check that it sits to the left of \textbf{\gu{શ}}; repair its placement once if needed.

\subsection*{Writing — recall after hiding}
Hide the model. Add short \emph{i} to \textbf{\gu{શ}} and write \emph{shi}. Uncover \textbf{\gu{શિ}}, compare, and stop after one repair.

\begin{tcolorbox}[breakable,colback=teal!4,colframe=teal!35!black,title={Before you move on}]
Is \textbf{\gu{િ}} seen before or after its consonant? (\textbf{Before.}) Is it pronounced before or after? (\textbf{After.})

Source: \href{https://www.unicode.org/charts/PDF/U0A80.pdf}{Unicode Gujarati chart}; placement and example follow \texttt{data/scripts/gujarati.json}.
\end{tcolorbox}
\section[retroflex lla — tha — independent a …]{\gu{ળ} — \gu{થ} — \gu{અ} — \gu{િ}: four forms return}

\label{lesson:GU-R04-first-four-r1}



\begin{quote}
Keep the four cards visible. Point as you say retroflex \emph{lla}, aspirated \emph{tha}, independent short \emph{a}, and attached short \emph{i}. For \textbf{\gu{િ}}, also say: seen before, heard after.
\end{quote}

\subsection*{Your turn}
Read the row once from left to right. Give \textbf{\gu{િ}} its role as an attached sign, not as an independent letter. Record a reading score out of four.

\subsection*{Writing — from sound or role}
Cover every model. Write the four cues in order. Record a separate writing score out of four; a reading point cannot replace a written form. Aim for \textbf{4/4} in each column, with one focused repair allowed.

\begin{tcolorbox}[breakable,colback=teal!4,colframe=teal!35!black,title={Before you move on}]
Write the only attached sign in the row, then say its placement rule.
\end{tcolorbox}
